\scanpage{40}
% L1–5
\margin{\unsure{$\vert\,\vert$|11}}
reste vraie pour le module $M/xM$ sur $A/xA$. Or ce dernier est encore Coh. Mac, et l'hyp. de récurrence s'applique. Donc \deleted{$M/xM$} $A/xA$ est Coh. Mac, donc $A$ l'est. D'autre part $M/xM$ \deleted{$\simeq\Omega_{A/xA}$ donc} est fond. pour $A/xA$, d'où

% L6–8
\[
\begin{aligned}
 D(M/xM)&=\text{annulateur de }x\text{ dans }D(M)\\
 \rotatebox{90}{$\simeq$}&\\
 H^{n}(A/xA)&=\text{annulateur de }x\text{ dans }H^n(A).
\end{aligned}
\]

% L9–18
Donc l'hom naturel [\unsure{s'ass à|déduit de} $H^n(M)\simeq H^n(A)\otimes M=D(A)$]
\[
 H^n(A)\longrightarrow D(M)
\]
est tel qu'il induit un isom sur les éléments annulés par $x$, donc (lemme) il est injectif, ou par transposition
\[
 M\to\Omega_A
\]
est surjectif. D'autre part il induit un isom pour les $H^n$, i.e. le $H^n$ du noyau est $0$, donc ce noyau est de dim $<n$. Comme $M$ est équidim, ce noyau est nul.

% L19–24
\deleted{Prenons maintenant (v), en supposant $M$ Coh. Mac.}

\deleted{Procédons} \inserted{Supposons} \deleted{$M=$ \unsure{socle|Spec}$(M)$} \deleted{\inserted{pour}} $M/x_1M\unsure{+\ldots|\ldots}+x_nM$ a un annulateur de $\mathfrak m$ de dim $1$, $M$ \deleted{Coh. Mac.} \inserted{Coh. Mac \uline{fidèle}.}

\deleted{\unsure{Prouvons}} Alors \inserted{(\unsure{§ 1|cf 1|cf. § 1}, prop 4, cor. 1),} $\operatorname{Ext}^{n}(\deleted{\unsure{A/|A/\underline{x}}}\underline{k},M)=\deleted{\operatorname{Hom}(A/\underline{x},\unsure{H^{n}(M)|M^{n}(M)|M/\underline{x}M})}\inserted{\operatorname{Hom}(k,M/\underline xM)\simeq k}$, donc on est ramené à \unsure{(i* e)|(i e)|(i et e)}

% archive sheet number: 39
